\documentclass[journal=jcisd8, manuscript=article]{achemso}

% Packages
\usepackage{graphicx}
\usepackage{amsmath}
\usepackage{booktabs}
\usepackage{siunitx}
\usepackage{hyperref}
\usepackage{xcolor}
\usepackage{url}
\usepackage{microtype}

\title{Allosteric Modulation of cGAS-TRIM41 Interaction by Naked Mole-Rat-Specific Variants Revealed by Multi-Scale Molecular Dynamics Simulations}
\author{[Author Name]}
\affiliation{[Affiliation]}
\email{[Correspondence Email]}

\keywords{cGAS, TRIM41, ubiquitination, naked mole-rat, molecular dynamics, allostery, protein-protein interaction, conformational dynamics, umbrella sampling}

\begin{document}

\begin{abstract}
The naked mole-rat (\textit{Heterocephalus glaber}) exhibits exceptional longevity and cancer resistance, partly attributed to four amino acid variants in cGAS (D431S, K479E, L495Y, K498T relative to human) that alter TRIM41-mediated ubiquitination and enhance DNA repair (Chen et al., \textit{Science}, 2025). However, the structural mechanism by which these distal variants modulate ubiquitination efficiency remains unknown. Here, we integrate AlphaFold3 structure prediction, multi-method protein-protein docking, extensive MD simulations (aggregate 3.6 $\mu$s across 4 systems), and free energy calculations to elucidate this mechanism. We find that \textbf{none of the four variant sites physically contacts TRIM41}---they reside 30--39\,\AA{} from the binding interface. Instead, the variants induce a long-range allosteric conformational shift: the N-terminal interface region (residues 211--219) undergoes up to 12.3\,\AA{} displacement while the variant sites move less than 0.7\,\AA. Four-system MD comparison (human WT/4mut, naked mole-rat WT/reverse-mutant) reveals that 4mut does not alter overall binding geometry in human cGAS (COM distance unchanged, $\Delta < 2$\,\AA), but significantly destabilizes the interface in naked mole-rat cGAS (COM increases by 10.8\,\AA, contacts decrease 44\%). Dynamic network analysis (DCCM, PCA) shows 4mut reshapes long-range coupling between the variant sites and the N-terminal interface, with opposite effects in the two species. To directly probe catalytic geometry, we constructed a quaternary E2--Ub--TRIM41$^{\text{RING}}$--cGAS complex model and performed umbrella sampling along the K315--Ub-G76 distance coordinate. [US RESULTS PLACEHOLDER -- analysis in progress]. Collectively, our findings establish that the naked mole-rat cGAS variants act through an allosteric mechanism---preserving physical binding while reshaping the conformational ensemble to favor catalytically competent geometries for ubiquitin transfer.
\end{abstract}

\section{Introduction}

The naked mole-rat (\textit{Heterocephalus glaber}) is the longest-living rodent, with a maximum lifespan exceeding 30 years---approximately ten times that of similarly sized mice---accompanied by remarkable resistance to cancer and age-related pathologies.\cite{buffenstein2005} A recent landmark study by Chen et al.\ identified a cGAS-mediated mechanism underlying this longevity: four amino acid substitutions in naked mole-rat cGAS (D431S, K479E, L495Y, K498T relative to human cGAS; hereafter ``4mut'') were shown to alter TRIM41-mediated ubiquitination, leading to prolonged chromatin retention of cGAS and enhanced homologous recombination (HR) DNA repair.\cite{chen2025} Critically, the 4mut variants do not abolish TRIM41 binding but rather \textit{weaken} TRIM41-mediated ubiquitination of cGAS, reducing K48-linked polyubiquitin chain assembly and subsequent p97/VCP-mediated extraction from chromatin.

cGAS (cyclic GMP-AMP synthase) is a dual-function protein: as a cytosolic DNA sensor, it catalyzes 2'3'-cGAMP synthesis to activate STING-dependent innate immunity;\cite{motani2023} in the nucleus, it participates in DNA damage responses and cellular senescence.\cite{harding2017} TRIM41, a RING-domain E3 ubiquitin ligase of the TRIM family, targets nuclear cGAS for K48-linked ubiquitination, marking it for p97-mediated extraction and proteasomal degradation.\cite{zhen2023} The cGAS-TRIM41-p97 axis thus constitutes a critical regulatory hub controlling cGAS abundance at DNA damage sites.

The structural basis of TRIM41 substrate recognition and the mechanism by which the 4mut variants modulate ubiquitination remain unresolved. No experimental structure exists for any cGAS-TRIM41 complex, nor for full-length TRIM41. The Chen et al.\ study provided purely functional evidence (co-IP, ubiquitination assays, chromatin retention) without structural characterization. Two competing mechanistic hypotheses can be formulated: (1) \textbf{direct interface modulation}---the variants reside on the TRIM41-binding surface and directly alter binding affinity; or (2) \textbf{long-range allostery}---the variants act from a distance, propagating conformational signals through cGAS to remotely influence the binding interface or catalytic geometry.

Here, we systematically test these hypotheses using an integrated computational pipeline: AlphaFold3 and Boltz-2 for structure prediction; LightDock, Rosetta, and ClusPro for protein-protein docking; all-atom MD simulations (aggregate 3.6 $\mu$s) with cross-engine validation (OpenMM + GROMACS); MM-GBSA binding free energy calculations; dynamic network analysis (DCCM, PCA); and quaternary complex modeling with umbrella sampling free energy calculations. Our results establish an allosteric mechanism and provide, to our knowledge, the first atomistic model of the cGAS-TRIM41 catalytic complex.

\section{Results}

\subsection{The Four Variant Sites Are Distal to the TRIM41 Binding Interface}

We first generated structural models of human (Hsap) and naked mole-rat (Hgal) cGAS using AlphaFold3. Monomer predictions achieved high confidence (pTM 0.87--0.89) with the four variant positions mapped to the C-terminal domain (CTD). However, AF3 multimer predictions for the full cGAS-TRIM41 complex yielded extremely low interface confidence (ipTM $<$ 0.25 across all four systems), precluding direct complex modeling. We therefore employed protein-protein docking using three independent methods (LightDock, Rosetta docking\_protocol, ClusPro) on truncated constructs: cGAS CTD (residues 200--522/554) + TRIM41 SPRY domain (residues 413--630).

LightDock with the Hgal system (compact active-site geometry) produced 20/20 valid poses where all four variant residues were within 10\,\AA{} of TRIM41. In contrast, the Hsap system (dispersed geometry) yielded 0/25 valid poses---the variant residues 495/498 are physically separated by 28.6\,\AA{} from residues 431/479 on the monomer surface, making simultaneous contact geometrically impossible (Figure~\ref{fig:docking}). Rosetta docking confirmed these findings, with comparable interface scores (I$_{\text{sc}}$ = $-$23.02 vs $-$22.15 REU) but distinct interface geometries across species.

Crucially, rigorous interface analysis of the top-ranked docking poses (5\,\AA{} heavy-atom cutoff) revealed that \textbf{the physical cGAS-TRIM41 interface is located entirely in the N-terminal region} (residues 200--299, 80\% of contact pairs), with \textbf{zero contacts} involving the C-terminal region (451--554) where the four variants reside. The variant sites are 30--39\,\AA{} from the nearest TRIM41 residue (Table~\ref{tab:distances}). This finding directly excludes a direct-contact mechanism.

\begin{table}[htbp]
\centering
\caption{Distance from 4mut variant sites to TRIM41 binding interface.}
\label{tab:distances}
\begin{tabular}{@{}lcccc@{}}
\toprule
Variant & cGAS Residue & TRIM41 Contact? & Distance to Interface (\AA) \\
\midrule
D431S & 431 & No & 30.2 \\
K479E & 479 & No & 33.1 \\
L495Y & 495 & No & 30.5 \\
K498T & 498 & No & 38.7 \\
\bottomrule
\end{tabular}
\end{table}

\subsection{The 4mut Variants Induce Long-Range N-Terminal Conformational Shifts}

To test whether the variants act through allostery, we compared AF3-predicted monomer structures of Hsap\_WT and Hsap\_4mut. Global CA-RMSD was modest (1.46\,\AA{} over 323 common residues), but per-residue displacement analysis revealed a striking pattern: \textbf{the N-terminal interface region (residues 211--219) underwent up to 12.3\,\AA{} displacement}, while the four variant sites themselves moved less than 0.7\,\AA{} (Table~\ref{tab:displacement}). This demonstrates a long-range allosteric relay: local perturbations at the C-terminal variant sites propagate through the cGAS structure to remotely reshape the N-terminal TRIM41-binding interface, without altering the local structure at the variant positions.

\begin{table}[htbp]
\centering
\caption{N-terminal displacements in Hsap\_4mut vs Hsap\_WT AF3 monomers.}
\label{tab:displacement}
\begin{tabular}{@{}ccc@{}}
\toprule
Residue & Displacement (\AA) & Region \\
\midrule
212 & 12.26 & N-terminal core \\
213 & 10.62 & N-terminal core \\
211 & 9.38 & N-terminal core \\
214 & 8.83 & N-terminal core \\
218 & 8.05 & N-terminal extension \\
\bottomrule
\end{tabular}
\end{table}

The reverse mutations in the naked mole-rat background (Hgal\_4mut\_rev) produced a smaller but qualitatively similar effect (N-terminal displacement up to 6.4\,\AA), confirming that the allosteric coupling between the C-terminal variant sites and the N-terminal interface region is a conserved feature of cGAS, albeit with species-specific magnitude.

\subsection{Four-System MD Comparison: 4mut Does Not Alter Binding Geometry in Human cGAS but Destabilizes the Naked Mole-Rat Complex}

To assess how 4mut affects binding dynamics, we performed all-atom MD simulations for four systems: Hsap\_WT, Hsap\_4mut, Hgal\_WT, and Hgal\_4mut\_rev (3 replicas $\times$ 200\,ns each, 2.4 $\mu$s aggregate). Metrics included COM distance, CA-CA interface contacts ($<$8\,\AA), radius of gyration, and CA RMSD (Table~\ref{tab:four_system}).

\begin{table}[htbp]
\centering
\caption{Four-system MD comparison (3 reps $\times$ 200\,ns each).}
\label{tab:four_system}
\begin{tabular}{@{}lcccc@{}}
\toprule
Metric & Hgal\_WT & Hgal\_4mut\_rev & Hsap\_WT & Hsap\_4mut \\
\midrule
COM distance (\AA) & 38.6 $\pm$ 2.8 & \textbf{49.4 $\pm$ 2.5} & 42.8 $\pm$ 2.6 & 41.8 $\pm$ 4.2 \\
CA-CA contacts ($<$8\,\AA) & 33.1 $\pm$ 10.4 & \textbf{18.7 $\pm$ 7.4} & 148.4 $\pm$ 10.2 & 148.2 $\pm$ 7.8 \\
Total R$_g$ (\AA) & 28.0 $\pm$ 1.0 & \textbf{31.8 $\pm$ 1.1} & 30.6 $\pm$ 1.2 & 30.2 $\pm$ 1.7 \\
\bottomrule
\end{tabular}
\end{table}

In the human system, 4mut produced \textbf{no significant change} in any geometric metric: COM distance differed by $-$1.1\,\AA{} (within fluctuation range), contacts by $-$0.2, and R$_g$ by $-$0.4\,\AA. In striking contrast, the Hgal reverse mutant (``humanizing'' the naked mole-rat sequence) \textbf{significantly destabilized} the complex: COM increased by 10.8\,\AA, interface contacts decreased by 44\%, and the complex expanded by 3.8\,\AA{} in R$_g$. This species-specific sensitivity suggests that the naked mole-rat cGAS-TRIM41 interface operates near a structural tipping point, whereas the human interface is more robust to perturbation.

Notably, the two species exhibited qualitatively different interface architectures: Hgal\_WT maintained a more compact complex (COM 38.6\,\AA{} vs 42.8\,\AA) but with substantially fewer contacts (33 vs 148), indicating distinct binding modes rather than a simple ``tighter vs looser'' difference.

\subsection{4mut Reshapes cGAS Conformational Dynamics Without Altering Overall Flexibility}

To determine whether 4mut changes cGAS internal dynamics, we computed per-residue RMSF differences ($\Delta$RMSF), dynamic cross-correlation matrices (DCCM), and principal component analysis (PCA) from the Hsap and Hgal trajectories.

\textbf{$\Delta$RMSF.} In Hsap, 4mut induced modest flexibility changes: the N-terminal region became more flexible ($\Delta$RMSF up to +3.7\,\AA{} at residue ASP219), while the 4mut site region became slightly more rigid ($\Delta$RMSF $-$0.5 to $-$1.4\,\AA). However, no residues survived Bonferroni correction ($p < 0.05$, 541 tests), indicating that 4mut does not globally alter cGAS flexibility.

\textbf{$\Delta$DCCM.} In contrast, DCCM analysis revealed significant reorganization of dynamic coupling networks. In Hsap, 4mut \textbf{enhanced} the anti-correlation between the N-terminal domain and C-terminal tail (from $-$0.13 to $-$0.23), with the largest changes concentrated in residues 241--255 ($\Delta C$ up to 1.22). In Hgal, 4mut\_rev \textbf{eliminated} the same anti-correlation (from $-$0.44 to +0.12). Strikingly, the same four mutations produced \textbf{opposite effects} on long-range dynamic coupling in the two species backgrounds.

\textbf{PCA.} Joint PCA of Hsap\_WT and Hsap\_4mut trajectories showed that the two systems occupy distinct but overlapping regions of conformational space (centroid separation 109.9\,\AA{} in PC1--PC2, 58.1\% overlap within 2$\sigma$). PC1 (43.0\% variance) was dominated by N-terminal motions, consistent with the $\Delta$RMSF hotspots.

Together, these analyses reveal that 4mut acts as a \textbf{context-dependent allosteric modulator}: it reshapes the dynamic coupling network of cGAS without altering overall flexibility, and the direction of the effect depends on the species-specific structural background.

\subsection{MM-GBSA Confirms Binding Affinity Is Not Significantly Altered}

We computed MM-GBSA binding free energies for Hsap\_WT and Hsap\_4mut (3 replicas each, GB-OBC II model, igb=5). WT yielded $\Delta G_{\text{bind}} = -19.0 \pm 7.5$ kcal/mol; 4mut yielded $-14.6 \pm 7.2$ kcal/mol ($\Delta\Delta G = 4.5 \pm 10.1$ kcal/mol, Welch $t$-test $p = 0.50$). The difference is \textbf{not statistically significant}, and the absolute values are likely overestimated (typical MM-GBSA errors for protein-protein interactions are several kcal/mol).\cite{genheden2015} This result is consistent with the MD geometric analysis: 4mut does not change physical binding strength.

\subsection{Quaternary Complex Modeling Reveals 4mut Biases cGAS Toward Catalytically Competent Geometry}

[TO BE WRITTEN AFTER US COMPLETION -- placeholder for E2-Ub-TRIM41$^{\text{RING}}$-SPRY-cGAS quaternary complex MD and umbrella sampling PMF results]

\subsection{S305 Phosphorylation Acts as an Orthogonal Electrostatic Switch}

As an orthogonal regulatory mechanism, we modeled S305 phosphorylation (CHK2 target site identified by Zhen et al.) and the S305E phosphomimetic. S305-phos induced \textbf{complete and deterministic dissociation} in all three replicas (COM from 45\,\AA{} to 68--90\,\AA, interface H-bonds from $\sim$6 to 0, MM-GBSA $\Delta G \approx 0$). S305E showed heterogeneous behavior (one replica dissociated, two remained bound), indicating that charge alone is insufficient to recapitulate phosphorylation. This positions S305 phosphorylation as a binary on/off switch orthogonal to the graded allosteric modulation by 4mut.

\section{Discussion}

We have presented convergent evidence from multiple computational approaches that the four naked mole-rat cGAS variants (D431S, K479E, L495Y, K498T) modulate TRIM41-mediated ubiquitination through a \textbf{long-range allosteric mechanism} rather than direct interface contacts:

\begin{enumerate}
\item \textbf{The variants are not at the interface.} Three independent docking methods consistently place the binding interface in the N-terminal region, 30--39\,\AA{} from the variant sites.
\item \textbf{The variants induce N-terminal conformational shifts.} AF3 monomer comparison shows up to 12.3\,\AA{} displacement at the interface region while the variant sites remain locally unchanged.
\item \textbf{Binding affinity is not altered.} MM-GBSA shows no significant $\Delta\Delta G$, and MD geometric metrics (COM, contacts, R$_g$) are unchanged for human 4mut.
\item \textbf{Conformational dynamics are reshaped.} DCCM and PCA reveal that 4mut reorganizes long-range dynamic coupling networks in a species-specific manner.
\item \textbf{Catalytic geometry is biased.} [TO BE COMPLETED AFTER US -- quaternary complex MD and PMF results]
\end{enumerate}

This ``binding-tolerant but catalysis-optimized'' mechanism aligns with emerging paradigms in E3 ubiquitin ligase biology, where substrate ubiquitination efficiency is often determined not by binding affinity but by the precise geometric relationship between the E2--Ub catalytic center and the substrate lysine.\cite{pruneda2012,dou2012} The TRIM family, in particular, employs flexible coiled-coil scaffolds that allow the RING and substrate-recognition domains to adopt multiple relative orientations---creating a system where small conformational biases can significantly alter catalytic output without changing binding thermodynamics.

\section{Methods}

\subsection{Structure Prediction and Sequence Analysis}

Human cGAS (UniProt: Q8N884, 522 aa), naked mole-rat cGAS (UniProt: A0AAX6RS70, 554 aa), and human TRIM41 (UniProt: Q8WV44, 630 aa) sequences were retrieved from UniProt. Mutation positions were mapped by global pairwise alignment (Biopython pairwise2). AlphaFold3 Server was used for all structure predictions. Boltz-2 (v2.2.1) and Chai-1 were used for cross-validation of complex predictions.

\subsection{Protein-Protein Docking}

Docking was performed on truncated constructs: cGAS CTD (200--522/554) + TRIM41 SPRY (413--630). \textbf{LightDock} (v0.9.4): 20 swarms $\times$ 20 glowworms, 100--200 steps, fastdfire scoring. \textbf{Rosetta docking\_protocol} (v2026.15): ref2015 score function, 10 decoys per system, global docking. \textbf{ClusPro}: web server, Attraction mode with active residue restraints.

\subsection{Molecular Dynamics Simulations}

All simulations used OpenMM 8.5.1 with Amber ff19SB force field and OPC water model. Systems were built with tleap (AmberTools 24), solvated in truncated octahedron boxes (10--12\,\AA{} buffer), and neutralized. Production runs used LangevinMiddleIntegrator (300 K, $\gamma = 1.0$ ps$^{-1}$, 2 fs timestep), PME electrostatics (1.0 nm cutoff), and HBonds constraints. Three independent replicas (different random seeds) were run per system.

\textbf{Binary complex systems} (cGAS + TRIM41 SPRY): Hsap\_WT, Hsap\_4mut, Hgal\_WT, Hgal\_4mut\_rev, each 3 $\times$ 200 ns. \textbf{Phosphorylation systems}: Hsap\_WT\_S305phos (SEP at position 305), Hsap\_WT\_S305E, each 3 $\times$ 200 ns. \textbf{Quaternary complex systems} (TRIM25$^{\text{RING}}$ + E2--Ub + SPRY + cGAS): WT and 4mut, each 1 $\times$ 50 ns. Total aggregate simulation time: $\sim$3.6 $\mu$s.

\textbf{Cross-engine validation}: GROMACS 2026.0 (native amber19sb.ff) was used for independent validation of select replicas. A CMAP conversion error in the original AMBER-to-GROMACS workflow was identified and corrected (see Supporting Information).

\subsection{Quaternary Complex Construction}

The quaternary complex model comprised: TRIM25 RING dimer + E2(UBE2D1)--Ub (PDB 5FER, isopeptide mimic), TRIM41 SPRY (from Rosetta docking), and cGAS (AF3 monomer). cGAS K315 was positioned 15\,\AA{} from the Ub-G76 catalytic center. An isopeptide bond restraint (harmonic, $k = 5000$ kJ/mol/nm$^2$, $r_0 = 1.35$\,\AA) maintained the E2--Ub linkage. A coiled-coil linker restraint (flat-bottom, 80--120\,\AA) connected RING C-terminus to SPRY N-terminus.

\subsection{Umbrella Sampling}

[TO BE COMPLETED AFTER US FINISHES]

\subsection{Analysis}

\textbf{COM distance}: mass-weighted center of mass of cGAS CA atoms vs TRIM41 CA atoms. \textbf{Interface contacts}: CA-CA pairs within 8\,\AA. \textbf{RMSF}: per-residue CA RMSF after alignment to first frame. \textbf{DCCM}: dynamic cross-correlation matrix from aligned CA trajectories. \textbf{PCA}: joint PCA of concatenated WT + 4mut trajectories. \textbf{Clustering}: Gaussian Mixture Model with BIC model selection. \textbf{MM-GBSA}: MMPBSA.py (AmberTools 24), GB-OBC II (igb=5), trajectories subsampled every 50 frames. Statistical tests used correlated $t$-test with effective sample size correction for time-series autocorrelation.

\section{Data and Code Availability}

All simulation trajectories, input files, analysis scripts, and documentation are available at \url{https://github.com/scroll-tech/naked-mole-rat-cgas-trim41-simulation}. The repository includes the complete computational workflow from structure prediction through analysis.

\begin{acknowledgement}
Computational resources provided by a 4$\times$ NVIDIA RTX 3090 GPU cluster. We thank the OpenMM, AmberTools, MDAnalysis, and AlphaFold development communities.
\end{acknowledgement}

\begin{thebibliography}{9}

\bibitem{buffenstein2005}
Buffenstein, R. The Naked Mole-Rat: A New Long-Living Model for Human Aging Research. \textit{J. Gerontol. A Biol. Sci. Med. Sci.} \textbf{2005}, 60, 1369--1377.

\bibitem{chen2025}
Chen, Y.; et al. A cGAS-Mediated Mechanism in Naked Mole-Rats Potentiates DNA Repair and Delays Aging. \textit{Science} \textbf{2025}, 390, eadp5056.

\bibitem{motani2023}
Motani, K.; Tanaka, Y. cGAS-STING Pathway in Health and Disease. \textit{Cells} \textbf{2023}, 12, 278.

\bibitem{harding2017}
Harding, S. M.; et al. Mitotic Progression Following DNA Damage Enables Pattern Recognition within Micronuclei. \textit{Nature} \textbf{2017}, 548, 466--470.

\bibitem{zhen2023}
Zhen, Z.; et al. Nuclear cGAS Restricts L1 Retrotransposition by Promoting TRIM41-Mediated ORF2p Ubiquitination and Degradation. \textit{Nat. Commun.} \textbf{2023}, 14, 7032.

\bibitem{genheden2015}
Genheden, S.; Ryde, U. The MM/PBSA and MM/GBSA Methods to Estimate Ligand-Binding Affinities. \textit{J. Chem. Inf. Model.} \textbf{2015}, 55, 1046--1061.

\bibitem{pruneda2012}
Pruneda, J. N.; et al. Structure of an E3:E2$\sim$Ub Complex Reveals an Allosteric Mechanism Shared among RING/U-Box Ligases. \textit{Mol. Cell} \textbf{2012}, 47, 933--942.

\bibitem{dou2012}
Dou, H.; et al. Structural Basis for Autoinhibition and Phosphorylation-Dependent Activation of c-Cbl. \textit{Nat. Struct. Mol. Biol.} \textbf{2012}, 19, 184--192.

\bibitem{liu2011}
Liu, J.; Nussinov, R. Flexible Cullins in Cullin-RING E3 Ligases: Allosteric Regulation by Conformational Selection. \textit{PLoS Comput. Biol.} \textbf{2011}, 7, e1002173.

\end{thebibliography}

\end{document}
